\chapter{A Concrete Naming Certificate for $\CompactModulusCoverUp$}
\label{ch:concrete-instances-compactmoduluscover-namecert}
\origin{ai}

The $\CompactModulusCoverUp$ interface names the finite source packet that sits between compactness, pointwise continuity, the $\FiniteNetMinimumFoldUp$ lower-bound fold, and the $\UniformModulusUp$ consumer. It records the compact finite-net probes and the pointwise-modulus ledgers before the rational fold chooses a single uniform precision, so the fold need not also carry the cover-formation certificate.

\begin{definition}[CompactModulusCover carrier]
\label{def:compact-modulus-cover-carrier}
A $\CompactModulusCoverUp$ carrier is a finite $\BHist$ packet
$$
K=(X,F,E,P,R,U,M,H,C,Q,N)
$$
with the following visible rows:
\begin{enumerate}
\item $X$ is the compact source row, read through $\CompactUp$ or $\CompactMetricUp$ as a finite-net source over the displayed $\MetricSpaceUp$ distance witnesses;
\item $F$ is the continuous map or function row, read through $\ContinuousMapUp$ or $\ContinuousFunctionUp$ with its pointwise-modulus ledger;
\item $E$ is the positive tolerance request supplied by the compact-continuous consumer;
\item $P$ is the nonempty finite $\mathsf{ProbeBundle}$ or compact-net row returned by the compact source certificate;
\item $R$ assigns to each bundled probe the source-center history, the covering-radius history, and the local distance witness that covers source histories by that center;
\item $U$ assigns to each bundled probe the pointwise-modulus precision row for the tolerance request $E$ at the corresponding center;
\item $M$ is the $\UniformModulusUp$-facing handoff row containing the same bundle, local radii, pointwise rows, and consumer routes, but no folded precision;
\item $H$ is the $\FiniteNetMinimumFoldUp$-facing radius family consisting of positive $\RatUp$ half-radius ledgers subordinate to the rows in $U$;
\item $C$ records componentwise $\hsame$ transports for the source row, continuous row, tolerance, bundle, covering rows, pointwise rows, radius family, and handoff row;
\item $Q$ records the displayed $\Cont$ routes for compact-net readback, pointwise-modulus readback, metric-distance use, rational-radius comparison, fold handoff, and uniform-modulus consumption;
\item $N$ is the combined provenance row over $\BHist$, $\hsame$, $\Cont$, $\Pkg$, and the local $\NameCert_{\CompactModulusCoverUp}$ certificate.
\end{enumerate}
The carrier has no coordinate for a host open cover, a selected subcover outside the finite compact-net row, a host-level Lebesgue number, an already folded uniform precision, or a real-completeness theorem. Its role is exactly the finite compact-net and pointwise-modulus source packet consumed by $\FiniteNetMinimumFoldUp$ and then by $\UniformModulusUp$.
\end{definition}
\leandef{BEDC.Derived.CompactModulusCoverUp.CompactModulusCoverCarrier}
\leanchecked{BEDC.Derived.CompactModulusCoverUp.CompactModulusCoverTasteGate\_single\_carrier\_alignment}

\begin{theorem}[CompactModulusCover NameCert obligations]
\label{thm:compact-modulus-cover-namecert-obligations}
The local $\NameCert_{\CompactModulusCoverUp}$ surface for the carrier of \autoref{def:compact-modulus-cover-carrier} consists exactly of carrier admission for
$$
K=(X,F,E,P,R,U,M,H,C,Q,N),
$$
classifier stability under the componentwise $\hsame$ rows in $C$, finite-net coverage exactness for the compact source rows in $P$ and $R$, pointwise-modulus compatibility for the continuous rows in $F$ and $U$, rational positive-radius compatibility for the $\RatUp$ family $H$, and consumer non-escape through the displayed $\Cont$, $\Pkg$, and local naming rows in $Q$ and $N$.
\end{theorem}
\begin{proof}
Project the displayed carrier. Its rows are the compact source, continuous map or function, tolerance request, finite probe bundle, covering rows, pointwise-modulus rows, uniform-modulus handoff row, positive rational radius family, transports, continuation routes, and local provenance.

The compact side is read only through the $\CompactUp$ or $\CompactMetricUp$ finite-net rows, together with the $\MetricSpaceUp$ distance witnesses attached to $R$. The continuous side is read only through the $\ContinuousMapUp$ or $\ContinuousFunctionUp$ pointwise-modulus rows stored in $U$.

The radius family $H$ is the only row passed to the $\FiniteNetMinimumFoldUp$ lower-bound fold, and the handoff row $M$ is the only row passed to the $\UniformModulusUp$ consumer before that fold has produced a precision. Componentwise $\hsame$ transports in $C$ and displayed $\Cont$ routes in $Q$ preserve those rows without introducing another selector.

Thus the named obligations exhaust the local certificate surface. Since no carrier coordinate stores a host open cover, a selected subcover, a host-level Lebesgue number, a folded uniform precision, or a real-completeness theorem, the certificate cannot export any of those objects.
\end{proof}

\begin{theorem}[CompactModulusCover finite-net source exhaustion]
\label{thm:compact-modulus-cover-finite-net-source-exhaustion}
For an accepted $\CompactModulusCoverUp$ carrier
$$
K=(X,F,E,P,R,U,M,H,C,Q,N),
$$
every source history used by the $\UniformModulusUp$ handoff is covered by a bundled center in the finite $\mathsf{ProbeBundle}$ row $P$, and every bundled center that contributes a pointwise-modulus row in $U$ contributes exactly one positive $\RatUp$ radius ledger to the $\FiniteNetMinimumFoldUp$-facing family $H$. The exhaustion certificate is internal to $P,R,U,H,C,Q,N$ and does not read an infinite cover or a host subcover choice.
\end{theorem}
\begin{proof}
Begin with the compact finite-net source row. The $\CompactUp$ or $\CompactMetricUp$ certificate supplies the nonempty finite bundle $P$ and the covering rows $R$; each covered source history is related to a bundled center through the displayed $\MetricSpaceUp$ distance witness.

For each bundled center, read the pointwise-modulus row in $U$ from the $\ContinuousMapUp$ or $\ContinuousFunctionUp$ ledger at the tolerance request $E$. The rational positive-radius row in $H$ is then the displayed half-radius ledger subordinate to that pointwise row.

The handoff row $M$ records exactly the same bundle and local rows for the $\UniformModulusUp$ consumer. The transports in $C$ and routes in $Q$ replay bundle membership, center coverage, pointwise-modulus readback, and rational-radius comparison componentwise.

Therefore the handoff consumes all and only the finite-net rows in $P$. No infinite cover, host-selected subcover, or additional compactness object is available from the carrier.
\end{proof}
\leanchecked{BEDC.Derived.CompactModulusCoverUp.CompactModulusCoverCarrier\_finite\_net\_source\_exhaustion}

\begin{theorem}[CompactModulusCover real-completion handoff]
\label{thm:compact-modulus-cover-real-completion-handoff}
Let a compact-continuous real consumer request a positive tolerance row $E$ for a $\ContinuousMapUp$ or $\ContinuousFunctionUp$ row over a compact $\MetricSpaceUp$ source. An accepted $\CompactModulusCoverUp$ carrier hands the consumer a finite $\UniformModulusUp$ source ledger and a $\FiniteNetMinimumFoldUp$ radius family, both over the same $\BHist$, $\hsame$, $\Cont$, $\Pkg$, and $\NameCert$ provenance. This handoff is sufficient for the later rational lower-bound fold and real-uniform-modulus read, and it does not claim Cauchy completion, host real completeness, or a global modulus chosen outside the packet.
\end{theorem}
\begin{proof}
Read the accepted carrier $K$. The compact side supplies the finite source coverage through $X,P,R$; the continuous side supplies pointwise-modulus rows through $F,U$; and the tolerance request is the displayed row $E$.

Apply \autoref{thm:compact-modulus-cover-finite-net-source-exhaustion}. It aligns the finite compact-net bundle with the pointwise rows and the positive rational radius family, so the $\FiniteNetMinimumFoldUp$ consumer receives precisely the family $H$ it needs for the lower-bound fold.

The $\UniformModulusUp$ consumer reads $M$ together with the same transport, route, provenance, and local naming rows $C,Q,N$. After the lower-bound fold has supplied a single precision, the consumer can route the usual compact-uniform-modulus argument through the bundled metric witnesses and pointwise rows.

No completion step is part of this packet. Real consumers may use the resulting uniform-modulus ledger downstream, but the $\CompactModulusCoverUp$ carrier itself exports only the finite compact-net and pointwise-modulus source packet.
\end{proof}

\begin{theorem}[CompactModulusCover obligation triad]
\label{thm:compact-modulus-cover-obligation-triad}
For an accepted $\CompactModulusCoverUp$ carrier
$$
K=(X,F,E,P,R,U,M,H,C,Q,N),
$$
the consumer surface before $\UniformModulusUp$ exposes exactly three obligations: compact finite-net coverage through $X,P,R$, pointwise-modulus compatibility through $F,E,U$, and rational lower-bound fold handoff through $H,M,C,Q,N$. It exports no host open cover, host-selected subcover, Lebesgue number, folded uniform precision, or real-completeness theorem.
\end{theorem}
\begin{proof}
First read the NameCert obligations. They identify carrier admission, classifier stability, finite-net coverage, pointwise-modulus compatibility, rational positive-radius compatibility, and consumer non-escape as the local certificate surface.

Next separate the three consumer obligations. Finite-net source exhaustion supplies the compact coverage route through $X,P,R$; the pointwise rows $F,E,U$ provide the local modulus compatibility; and the real-completion handoff sends the radius family and structural rows through $H,M,C,Q,N$ toward the lower-bound fold.

The non-escape clauses in those theorems rule out every unlisted analytic object. Since the carrier contains no host open cover, host-selected subcover, Lebesgue number, folded uniform precision, or real-completeness theorem, the triad is the whole pre-$\UniformModulusUp$ surface.
\end{proof}

\begin{theorem}[CompactModulusCover AI TasteGate scope]
\label{thm:compact-modulus-cover-ai-tastegate-scope}
Any obligation-closure read for $\CompactModulusCoverUp$ must pass through the same single-carrier packet
$$
K=(X,F,E,P,R,U,M,H,C,Q,N)
$$
of \autoref{def:compact-modulus-cover-carrier}. The read must include the finite-net source exhaustion of \autoref{thm:compact-modulus-cover-finite-net-source-exhaustion} and must preserve carrier alignment for the finite-net rows, pointwise-modulus rows, fold handoff, transport rows, route rows, and local $\NameCert$ row before any closure promotion is recorded.
\end{theorem}
\begin{proof}
Project the carrier $K$. The compact finite-net source is read from $X,P,R$, and the pointwise-modulus data are read from $F,E,U$.

Apply finite-net source exhaustion. It ties every source history used by the $\UniformModulusUp$ handoff to the finite bundle $P$, and it ties every bundled center with a pointwise row in $U$ to the positive rational radius family $H$.

The fold-facing and consumer-facing rows are then fixed. The row $H$ is the $\FiniteNetMinimumFoldUp$ radius family, and the row $M$ is the $\UniformModulusUp$ handoff row before a folded precision is available.

Finally read the structural rows $C,Q,N$. They provide componentwise transport, displayed routes, and the local naming surface for the same packet, so an obligation-closure read cannot bypass the displayed carrier or use a second source packet.
\end{proof}

\closureat{CompactModulusCoverUp}{seedStr}

\begin{closurestatus}{\CompactModulusCoverUp}
  \constructivestory{A $\CompactModulusCoverUp$ packet is generated from a compact source row, a continuous map or function row, a positive tolerance request, a finite $\mathsf{ProbeBundle}$, compact-net coverage rows, pointwise-modulus rows, positive $\RatUp$ radius ledgers, componentwise $\hsame$ transports, displayed $\Cont$ routes, package provenance, and a local $\NameCert$ row.}
  \theoryclosure{\seedClosure}
  \scopeclosed{\autoref{def:compact-modulus-cover-carrier}, \autoref{thm:compact-modulus-cover-namecert-obligations}, \autoref{thm:compact-modulus-cover-finite-net-source-exhaustion}, and \autoref{thm:compact-modulus-cover-real-completion-handoff} register the finite compact-net and pointwise-modulus packet used before the lower-bound fold and uniform-modulus handoff.}
  \formalstatus{\formalTargetV}
  \bridgestatus{none}
  \origin{ai}
  \notclaimed{This chapter does not close the rational lower-bound fold, the final uniform-modulus theorem, Cauchy completion, host real completeness, external subcover selection, or bridge checking.}
  \upgradepath{Scoped closure requires Lean encoding and checked carrier admission, classifier stability, finite-net coverage exactness, pointwise-modulus compatibility, positive-radius compatibility, and consumer non-escape for the same displayed packet.}
\end{closurestatus}
